For two nonzero numbers $u<v$, $\dfrac{1}{u}-\dfrac{1}{v}=\dfrac{v-u}{uv}$. If they have the same sign, then $uv>0$, so the difference is positive. Thus in either same-sign case, $\dfrac{1}{a}>\dfrac{1}{b}>\dfrac{1}{c}$. If the signs differ, the denominator may be negative. For example, $-2<-1<1$, but their reciprocals are $-\dfrac{1}{2},-1,1$, which do not form a decreasing chain.
